\chapter*{Conclusion}
\addcontentsline{toc}{chapter}{Conclusion}

This thesis presented DBMS-Nautilus, a grammar-based greybox fuzzing system for automated vulnerability detection in SQLite. What makes DBMS-Nautilus different is its grammar: the grammar determines what the fuzzer can discover, while the fuzzing engine navigates within that space but cannot expand it. The central contribution is a grammar engineering methodology based on structural primitives --- SQL patterns extracted from CVE root-cause analysis, encoded as composable production rules without hardcoding any proof-of-concept inputs. The cross-pollination property of this design, where a single non-terminal serves multiple CVEs simultaneously and the fuzzer freely composes patterns across independent non-terminal groups, enables compositional discovery of vulnerability-triggering inputs.

The grammar contains 526 production rules, iteratively refined based on experimental evidence.

Regarding the first research question, DBMS-Nautilus successfully rediscovered all four target CVEs through compositional grammar generation: CVE-2020-13434 (integer overflow in printf, 14/15 version-runs), CVE-2020-13435 (null pointer dereference, 5/5 runs), CVE-2020-13871 (use-after-free in window function processing, 4/5 runs), and CVE-2020-15358 (heap buffer overflow, 5/5 runs). None of these were found by embedding proof-of-concept inputs --- the fuzzer independently composed the triggering SQL from separate structural primitives. The EBNF-Baseline found only 1 of 4 CVEs (CVE-2020-13871), confirming that 3 of the 4 target vulnerabilities require domain-specific structural patterns.

Regarding the second research question, DBMS-Nautilus discovered 13 distinct bug classes across 48 unique stack hashes spanning four SQLite versions: 4 are CVE rediscoveries without dedicated seed rules (BC001=CVE-2020-13435, BC003=CVE-2020-13434, BC006=CVE-2020-9327, BC007=CVE-2020-13871), 6 are non-CVE bugs that cause crashes on production builds across up to 10 consecutive releases, and 3 are sanitizer-only detections. Nine of the 13 bug classes were found exclusively by DBMS-Nautilus. The discovery of bugs beyond the grammar's explicit design targets confirms the cross-pollination hypothesis: structural patterns intended for specific CVEs also exercise adjacent code paths containing previously undocumented defects.

Regarding the third research question, the proposed grammar sacrifices throughput ($2.1\times$ slower) and edge coverage (14.1\% lower) in exchange for $108\times$ more unique root causes per campaign compared to a generic EBNF baseline grammar, with large effect sizes ($d = 1.00$, $p < 0.05$) on all four versions. This asymmetry is consistent with the hypothesis that structural complexity concentrates fuzzing effort on vulnerability-relevant code paths at the cost of raw execution speed.

Several limitations bound these findings. All experiments target a single DBMS (SQLite), and the results may not generalize to larger server-based systems. The context-free grammar cannot enforce semantic constraints, meaning a fraction of generated inputs are rejected at runtime. The sanitizer-based oracle misses non-crash bugs such as infinite loops and incorrect query results. The grammar design process requires manual domain expertise in both SQL internals and vulnerability patterns.

Four directions warrant further investigation. First, reinforcement learning integration could enable adaptive weight tuning based on coverage and crash-frequency signals, potentially accelerating discovery of bugs requiring rare compositional combinations. Second, cross-DBMS grammar portability should be tested by adapting the structural primitives methodology to MySQL and PostgreSQL. Third, semantic-aware generation extending the grammar with type constraints would reduce wasted executions on semantically invalid inputs. Fourth, automated grammar expansion using large language model analysis of vulnerability disclosures could partially automate the grammar design process.
